\documentclass[UTF8,zihao=-4]{ctexart}
\usepackage[a4paper,top=2.15cm,bottom=2.0cm,left=2.5cm,right=2.5cm]{geometry}
\usepackage{amsmath,amssymb,booktabs,array,enumitem,xcolor,fancyhdr}
\usepackage[hidelinks]{hyperref}
\setCJKmainfont[BoldFont=SimHei]{SimSun}
\setmainfont{Times New Roman}
\definecolor{inkblue}{HTML}{17456D}
\definecolor{softgray}{HTML}{555F68}
\setlength{\parindent}{2em}
\setlength{\parskip}{0.35em}
\setlist[itemize]{leftmargin=2em,itemsep=0.4em,topsep=0.25em}
\renewcommand{\baselinestretch}{1.28}
\pagestyle{fancy}
\fancyhf{}
\fancyhead[L]{\small\color{softgray}毕业论文开题思路}
\fancyhead[R]{\small\color{softgray}资产配置与最优化方法}
\fancyfoot[C]{\small\thepage}
\renewcommand{\headrulewidth}{0.4pt}
\setlength{\headheight}{15pt}
\begin{document}

\begin{center}
{\zihao{2}\bfseries\color{inkblue}基于宽松风险平价的\\[0.25em]多资产 ETF 配置研究}\\[0.8em]
{\large 本科毕业论文开题思路（拟）}\\[0.5em]
{\normalsize 研究方向\quad 资产配置与最优化方法}
\end{center}
\vspace{0.4em}

\section*{一、研究缘起与核心问题}

\subsection*{从实习中的关注出发}

实习期间接触到桥水全天候策略的风险分散理念，由此对多资产配置产生了兴趣。不同资产会受到不同市场因素的影响，如何将它们组合起来、避免组合过度依赖某一类风险，是希望进一步研究的问题。

结合老师最优化理论与应用的研究方向，拟以境内上市 ETF 为对象，把资产配置中的取舍写成一个明确的优化问题：在保持风险分散的基础上，如何考虑收益目标并确定资产权重？

\subsection*{希望回答的问题}

风险平价从各资产承担的风险出发进行配置，而收益预测本身又存在误差。若为了收益目标调整权重，怎样兼顾风险分散，并避免配置对预测过于敏感？

论文拟以风险预算组合作为参考，允许最终配置在一定程度上偏离参考权重，研究这种调整带来的收益与风险权衡。初步考虑多头、满仓和无杠杆条件，使研究范围集中在资产之间的资金分配。

\subsection*{选题的主要落点}

桥水全天候的配置思想提供研究动机，论文拟落在连续权重的最优化建模上。重点是说明模型如何表达投资取舍、在什么条件下能够求解，以及输入变化时配置是否稳定。实证比较用于解释和检验模型设计。

\clearpage
\section*{二、模型思路与最优化方法}

\subsection*{以风险预算组合作为起点}

初步考虑先通过风险预算模型得到参考权重，再围绕参考权重进行调整。这里的“宽松”是允许最终配置适度偏离参考组合，以容纳收益目标，而不是要求每一步都严格满足风险平价。

调整时拟考虑三类代价：偏离参考权重、组合波动，以及预测收益未达到目标。收益目标可参考基准组合的预测收益。希望先确定这些代价怎样表达、哪些约束必须保留，再讨论参数设置。

\subsection*{为什么考虑凸优化}

拟选择具有凸形式的风险预算模型，并将后续权重调整构造为凸优化问题。例如，协方差矩阵半正定时，组合方差是凸二次函数；权重非负、权重之和为一是线性约束。再采用凸的偏离与收益不足惩罚，可形成便于分析和求解的模型。

数学部分拟围绕三个问题展开：可行解和最优解是否存在；在什么条件下最优权重唯一；如何用最优性条件解释权重调整。具体证明深度希望由老师指导确定。

\subsection*{与老师研究方向的衔接}

老师主页列出的方向为最优化理论与应用、博弈论与应用，相关成果涉及鲁棒优化与变分不等式。\footnote{研究方向依据：\href{https://math.swufe.edu.cn/info/1091/3222.htm}{西南财经大学数学学院·王磊老师主页}。}据此，拟先以凸优化建模和解的性质作为衔接点；若范围合适，再讨论收益预测存在误差时，能否通过鲁棒优化改善配置稳定性。该延伸是否纳入论文，留待开题讨论。

\clearpage
\section*{三、研究安排与希望讨论的问题}

\subsection*{先完成基础模型}

首先梳理风险预算与凸优化配置的相关文献，明确参考组合与最终权重的关系。随后确定一个规模适当的基础模型，说明目标函数的含义、约束条件与凸性，并讨论最优解的性质。

理论分析拟围绕所选模型展开，以已有最优化方法为基础。开题阶段希望先明确应证明哪些性质、分析应达到什么深度，再确定求解与实证方案。

\subsection*{数据与比较思路}

数据拟选取境内上市、覆盖股票、债券、商品及海外市场的 ETF。具体资产范围和样本区间结合数据可得性与研究问题确定，并考虑基金上市时间不同带来的影响。

实证拟以风险预算参考组合为主要对照，并选取等权、固定股债比例等简单基准。除收益、波动与回撤外，重点观察收益预测小幅变化时，权重和交易需求是否明显改变。敏感性检验可先帮助识别问题，再判断是否需要引入鲁棒模型。

\subsection*{希望请老师指导的三个问题}

\begin{itemize}
\item 以风险预算为参考、兼顾收益目标的凸优化模型，是否适合作为本科论文的主线？
\item 数学部分宜重点讨论凸性、解的存在与唯一性，还是进一步分析最优性条件？
\item 对预测误差，先做敏感性检验是否足够，还是有必要增加一个简单的鲁棒优化扩展？
\end{itemize}

预期先形成“模型构建与性质分析—数值求解—实证解释”的基本结构，具体范围根据指导意见收敛。

\end{document}
